% -----------------------------------------------------------
\section{Humility, Humble}
% -----------------------------------------------------------
	\label{HUMILITY}
	
% ----------------------
\subsection{See also}
% ----------------------

	\hyperref[LOWLINESS]{Lowliness},

% ----------------------
\subsection{1828 American Dictionary of the English Language} % *1828
% ----------------------
	\label{HUMILITY-1828}
	
	% -----
	\subsubsection{Humility}
	% -----

		\begin{compactitem}
			\item[1] In ethics, freedom from pride and arrogance; humbleness of mind; a modest estimate of one's own worth. In theology, \textit{humility} consists in lowliness of mind; a deep sense of one's own unworthiness in the sight of God, self-abasement, penitence for sin, and submission to the divine will.
			\item[2] Act of submission.
		\end{compactitem}
		
	% -----
	\subsubsection{Humbleness}
	% -----

		\begin{compactitem}
			\item The state of being humble or low; humility; meekness.
		\end{compactitem}
		
	% -----
	\subsubsection{Humble (adjective)}
	% -----

		\begin{compactitem}
			\item[1] Low; opposed to high or lofty.
			\item[2] Low; opposed to lofty or great; mean; not magnificent; as a humble cottage.
			\item[3] Lowly; modest; meek; submissive; opposed to proud, haughty, arrogant or assuming. In an evangelical sense, having a low opinion of one's self, and a deep sense of unworthiness in the sight of God.
		\end{compactitem}
		
	% -----
	\subsubsection{Humble (verb)}
	% -----

		\begin{compactitem}
			\item To abase; to reduce to a low state.
			\item[1] To crush; to break; to subdue.
			\item[2] To mortify.
			\item[3] To make humble or lowly in mind; to abase the pride of; to reduce arrogance and self-dependence; to give a low opinion of one's moral worth; to make meek and submissive to the divine will; the evangelical sense.
			\item[4] To make to condescend.
			\item[5] To bring down; to lower; to reduce.
			\item[6] To deprive of chastity.
			\item To humble one's self, to repent; to afflict one's self for sin; to make contrite.
		\end{compactitem}

% % ----------------------
% \subsection{Oxford English Dictionary} % *OED
% % ----------------------
	% \label{HUMILITY-OED}

	% \begin{compactitem}
		% \item[]
	% \end{compactitem}

% ----------------------
\subsection{Scriptural Usage} % *SCRIPTURES
% ----------------------
	\label{HUMILITY-SCRIPTURES}

	% % -----
	% \subsubsection{Old Testament} % *OT
	% % -----
		% \label{HUMILITY-OT}
	
		% % --
		% \paragraph{KJV}
		% % --
			% \label{HUMILITY-OT-KJV}
	
			% \begin{tabular}{ll}
				% Total occurrences in the OT & 
			% \end{tabular}
			
			% \begin{compactdesc}
				% \item[]
			% \end{compactdesc}
			
		% % --
		% \paragraph{Hebrew Bible}
		% % --
			% \label{HUMILITY-HEBREW}
		
			% %
			% \subparagraph{\texorpdfstring{\he{sample word}}{}}
			% %
				% \label{PAGA} % whatever is in step bible without periods
			
				% \begin{compactdesc}
					% \item[HALOT , PDF ] \textbf{}
						% \begin{compactitem}
							% \item[1]
						% \end{compactitem}
					% \item[BDB , PDF ] \textbf{}
						% \begin{compactitem}
							% \item[1]
						% \end{compactitem}
					% \item[DCH PDF ] \textbf{}
						% \begin{compactitem}
							% \item[1]
						% \end{compactitem}
				% \end{compactdesc}
				
				% \begin{compactdesc}
					% \item[Some OT uses] \textbf{}
						% \begin{compactdesc}
							% \item[]
						% \end{compactdesc}
				% \end{compactdesc}
			
		% % --
		% \paragraph{Septuagint}
		% % --
			% \label{HUMILITY-LXX}
		
			% %
			% \subparagraph{\texorpdfstring{\gr{sample word}}{}}
			% %
		
	% -----
	\subsubsection{New Testament} % *NT
	% -----
		\label{HUMILITY-NT}
	
		% --
		\paragraph{KJV}
		% --
			\label{HUMILITY-NT-KJV}
			
	
			\begin{tabular}{ll}
				Total occurrences in the NT & 15
			\end{tabular}
			
			\begin{compactdesc}
				\item[Matthew 18:4] Whosoever therefore shall humble himself as this little child, the same is greatest in the kingdom of heaven.
				\item[Matthew 23:12] And whosoever shall exalt himself shall be abased; and he that shall humble himself shall be exalted.
				\item[Acts 20:19] Serving the Lord with all humility of mind, and with many tears, and temptations, which befell me by the lying in wait of the Jews:
				\item[Philippians 2:8] And being found in fashion as a man, he humbled himself, and became obedient unto death, even the death of the cross. \textit{(Note that this verse is speaking of Christ.)}
				\item[Colossians 2:18] Let no man beguile you of your reward in a voluntary humility and worshipping of angels, intruding into those things which he hath not seen, vainly puffed up by his fleshly mind,
				\item[Colossians 3:12] Put on therefore, as the elect of God, holy and beloved, bowels of mercies, kindness, humbleness of mind, meekness, longsuffering;
				\item[1 Peter 5:6] Humble yourselves therefore under the mighty hand of God, that he may exalt you in due time:
			\end{compactdesc}
			
		% --
		\paragraph{Greek New Testament} % *GREEK
		% --
			\label{HUMILITY-GREEK}
		
			%
			\subparagraph{\texorpdfstring{\gr{ταπεινόω}}{}}
			%
				\label{TAPEINOO}
			
				\begin{compactdesc}
					\item[BDAG 880a] \textbf{}
						\begin{compactitem}
							\item[1] to cause to be at a lower point, \textit{lower}
							\item[2] to cause someone to lose prestige or status, \textit{humble, humiliate, abase}
							\item[2a] with focus on reversal of status \gr{ταπ. ἑαυτὸν} \textit{humble oneself}
							\item[2b] with focus on shaming, with accusative of person or thing treated in this manner
							\item[2c] with focus on punitive aspect
							\item[3] to cause to be or become humble in attitude, \textit{humble, make humble}
							\item[4] to subject to strict discipline, \textit{constrain, mortify}
						\end{compactitem}
					\item[LSJ 1757a, PDF 1828] \textbf{}
						\begin{compactitem}
							\item[I] \textit{lower, reduce, decrease}
							\item[II.1] \textit{lessen}
							\item[II.2a] \textit{humble, abase}
							\item[II.2b] \textit{violate} a woman
							\item[II.3] in moral sense, \textit{make lowly, humble}
								\begin{compactitem}
									\item Mentions verses in Matthew and 1 Peter, the latter of which is a passive; according to the LSJ, the passive can mean ``humble oneself,'' but there are many instances in the NT where you get the word with the reflexive pronoun.
								\end{compactitem}
							\item[II.4] especially of fasting or abstinence
						\end{compactitem}
				\end{compactdesc}
				
				\begin{compactdesc}
					\item[Some NT uses] \textbf{}
						\begin{compactdesc}
							\item[{\hyperref[MATTHEW-18-4]{Matthew 18:4}}]
							\item[{\hyperref[MATTHEW-23-12]{Matthew 23:12}}]
							\item[{\hyperref[PHILIPPIANS-2-8]{Philippians 2:8}}]
							\item[{\hyperref[1PETER-5-6]{1 Peter 5:6}}]
						\end{compactdesc}
				\end{compactdesc}
				
			%
			\subparagraph{\texorpdfstring{\gr{ταπεινοφροσύνη}}{}}
			%
				\label{TAPEINOPHROSUNE}
			
				\begin{compactdesc}
					\item[BDAG 880a] \textbf{}
						\begin{compactitem}
							\item \textit{humility, modesty}
						\end{compactitem}
					\item[LSJ 1757a, PDF 1828] \textbf{}
						\begin{compactitem}
							\item \textit{humility} \textit{(The example given here is Ephesians 4:2.)}
						\end{compactitem}
				\end{compactdesc}
				
				\begin{compactdesc}
					\item[Some NT uses] \textbf{}
						\begin{compactdesc}
							\item[Acts 20:19]
							\item[Ephesians 4:2] \gr{μετὰ πάσης ταπεινοφροσύνης καὶ πραΰτητος, μετὰ μακροθυμίας, ἀνεχόμενοι ἀλλήλων ἐν ἀγάπῃ,} \textit{(This is translated as ``lowliness'' in the KJV.)}
							\item[Colossians 2:18] 
							\item[Colossians 3:12] 
							\item[Philippians 2:3] \gr{μηδὲν κατ᾽ ἐριθείαν μηδὲ κατὰ κενοδοξίαν, ἀλλὰ τῇ ταπεινοφροσύνῃ ἀλλήλους ἡγούμενοι ὑπερέχοντας ἑαυτῶν,}
						\end{compactdesc}
				\end{compactdesc}
				
			%
			\subparagraph{\texorpdfstring{\gr{ταπεινός}}{}}
			%
				\label{TAPEINOS}
			
				\begin{compactdesc}
					\item[BDAG 879b] \textbf{}
						\begin{compactitem}
							\item[1] pertaining to being of low social status or to relative inability to cope, \textit{lowly, undistinguished, of no account}
							\item[2] pertaining to being servile in manner, \textit{pliant, subservient, abject} a negative quality that would make one lose face in the Greco-Roman world, opposite of a free person's demeanor
							\item[3] pertaining to being unpretentious, \textit{humble}
						\end{compactitem}
					\item[LSJ 1756b, PDF 1827] \textbf{}
						\begin{compactitem}
							\item[1] of place, \textit{low-lying}
							\item[2] of persons, \textit{humbled, abased in power, pride}
							\item[3] of the spirits, \textit{downcast, dejected}
							\item[4] in moral sense, either bad, \textit{mean, base, abject}, or good, \textit{lowly, humble} \textit{(Here it also mentions that this meaning occurs frequently in the NT, and gives Matthew 11:29 and 2 Corinthians 7:6 as examples.)}
							\item[5] of things, \textit{mean, low, poor}
						\end{compactitem}
				\end{compactdesc}
				
				\begin{compactdesc}
					\item[Some NT uses] \textbf{}
						\begin{compactdesc}
							\item[Matthew 11:29] \gr{ἄρατε τὸν ζυγόν μου ἐφ᾽ ὑμᾶς καὶ μάθετε ἀπ᾽ ἐμοῦ, ὅτι πραΰς εἰμι καὶ ταπεινὸς τῇ καρδίᾳ, καὶ εὑρήσετε ἀνάπαυσιν ταῖς ψυχαῖς ὑμῶν·}
							\item[2 Corinthians 7:6] \gr{ἀλλ᾽ ὁ παρακαλῶν τοὺς ταπεινοὺς παρεκάλεσεν ἡμᾶς ὁ θεὸς ἐν τῇ παρουσίᾳ Τίτου·}
						\end{compactdesc}
				\end{compactdesc}
		
	% -----
	\subsubsection{Book of Mormon} % *BOM
	% -----
		\label{HUMILITY-BOM}
	
		% \begin{tabular}{ll}
			% Total occurrences in the BOM & 
		% \end{tabular}
		
		\begin{compactdesc}
			\item[{\hyperref[MOSIAH-21-13]{Mosiah 21:13--17}}]
			\item[{\hyperref[HELAMAN-3-35]{Helaman 3:35}}]
			\item[{\hyperref[3NEPHI-6-13]{3 Nephi 6:13}}]
		\end{compactdesc}
		
	% -----
	\subsubsection{Doctrine and Covenants} % *DC
	% -----
		\label{HUMILITY-DC}
	
		% \begin{tabular}{ll}
			% Total occurrences in the DC & 
		% \end{tabular}
		
		\begin{compactdesc}
			\item[{\hyperref[DC63-55]{DC 63:55}}]
		\end{compactdesc}
		
	% % -----
	% \subsubsection{Pearl of Great Price} % *PGP
	% % -----
		% \label{HUMILITY-PGP}
	
		% \begin{tabular}{ll}
			% Total occurrences in the PGP & 
		% \end{tabular}
		
		% \begin{compactdesc}
			% \item[]
		% \end{compactdesc}
		
% ----------------------
\subsection{Key Phrases} % *KEYPHRASES *PHRASES
% ----------------------
	\label{HUMILITY-PHRASES}

	\begin{compactdesc}
	
		\item[truly humble] \textbf{6} | \hyperref[ALMA-27-18]{Alma 27:18}; \hyperref[ALMA-32-6]{Alma 32:6}; \hyperref[ALMA-32-14]{Alma 32:14}; \hyperref[ALMA-32-15]{Alma 32:15}; \hyperref[DC54-3]{DC 54:3}; \hyperref[DC97-1]{DC 97:1}
			\begin{compactdesc}
				\item[Bible anchor verses] ---
					% \begin{compactdesc}
						% \item[Romans 5:5] \gr{}
					% \end{compactdesc}
				\item[Commentary] The phrase ``true humility'' doesn't appear at all.
			\end{compactdesc}
			
		% \item[] \textbf{\#times} | list
			% % \begin{compactdesc}
				% % \item[Bible anchor verses] \phantom{.}
					% % \begin{compactdesc}
						% % \item[Romans 5:5] \gr{}
					% % \end{compactdesc}
				% % \item[Commentary] ---
			% % \end{compactdesc}
			
	\end{compactdesc}
	
% % ----------------------
% \subsection{Bible Dictionaries} % *DICTIONARIES
% % ----------------------
	% \label{HUMILITY-BD}

% % ----------------------
% \subsection{Restoration Usage} % *RESTORATION
% % ----------------------
	% \label{HUMILITY-RESTORATION}

	% % -----
	% \subsubsection{Joseph Smith} % *JS
	% % -----
		% \label{HUMILITY-JS}
	
		% \begin{compactdesc}
			% \item[DATE/SOURCE]
		% \end{compactdesc}
	
	% % -----
	% \subsubsection{Brigham Young} % *BY
	% % -----
		% \label{HUMILITY-BY}
	
		% \begin{compactdesc}
			% \item[DATE/SOURCE]
		% \end{compactdesc}
	
% ----------------------
\subsection{Commentary and Studies} % *COMMENTARY *STUDIES
% ----------------------
	\label{HUMILITY-COMMENTARY}

	% % -----
	% \subsubsection{Key Points}
	% % -----
		% \label{HUMILITY-KEYS}
	
		% \begin{compactitem}
			% \item
		% \end{compactitem}
		
	% -----
	\subsubsection{Humbling yourself versus being compelled to be humble}
	% -----
		\label{HUMILITY-yourself}
		
		\hyperref[ALMA-32-14]{Alma 32:14} and the scriptures around there are key for this. See also:
		
			\begin{compactitem}
				\item \hyperref[DC61-37]{DC 61:37}
			\end{compactitem}